\documentclass[11pt,a4paper]{article}

\usepackage{amsmath,amssymb}

\title{SoICT 2026: Integrated Passenger and Parcel Transportation}
\date{}

\begin{document}
\maketitle

\section{Input}

Let $\mathcal K=\{1,\ldots,K\}$, $\mathcal N=\{1,\ldots,N\}$, and
$\mathcal M=\{1,\ldots,M\}$ index taxis, passenger requests, and parcel
requests, respectively. The input consists of:

\begin{center}
\renewcommand{\arraystretch}{1.35}
\begin{tabular}{@{}p{0.23\linewidth}p{0.70\linewidth}@{}}
\hline
\textbf{Input} & \textbf{Given parameters} \\
\hline
Taxi $k\in\mathcal K$
    & $O_k$: initial location; $Q_k$: parcel capacity. Available at time $0$. \\
\hline
Passenger $i\in\mathcal N$
    & $P_i^n,D_i^n$: pickup and drop-off points; $T_i^n$: revenue;
      $[E_i^n,L_i^n]$: pickup time window; $S_i^n$: service duration at
      each endpoint. \\
\hline
Parcel $j\in\mathcal M$
    & $P_j^h,D_j^h$: pickup and drop-off points; $T_j^h$: revenue;
      $w_j$: weight; $[E_j^{h,p},L_j^{h,p}]$: pickup time window;
      $[E_j^{h,d},L_j^{h,d}]$: delivery time window; $S_j^h$: service
      duration at each endpoint. \\
\hline
Travel $(u,v)$
    & $t(u,v)\geq 0$: travel time; $c(u,v)\geq 0$: travel cost for every
      ordered pair of points. Neither function need be symmetric. \\
\hline
\end{tabular}
\end{center}

The network has $K+2N+2M$ labeled vertices: one starting point per taxi and
two service nodes per request.

\section{Output}

A schedule consists of one route per taxi,
\[
    R_k=(v_0^k,v_1^k,\ldots,v_{m_k}^k),\qquad
    v_0^k=O_k,\quad m_k\geq 0,
\]
where subsequent nodes are the pickup and drop-off points of served requests.
The routes identify which requests are served and which taxi serves each one.
An unused taxi has $m_k=0$. Service times must make every route feasible under
the constraints below. The probation output begins with the benefit, followed by exactly $K$ lines.
Line $k$ contains $m_k$ and the $m_k$ service vertices, using the input's
one-based labels and omitting the origin; an unused taxi has a line containing $0$.

\section{Objective}

Let $x_i^n,x_j^h\in\{0,1\}$ indicate whether passenger $i$ and parcel $j$ are
served. Maximize total revenue from served requests minus the travel cost of all
taxi routes:
\[
\begin{aligned}
    \max_{\text{feasible }R_1,\ldots,R_K}\quad
    &\sum_{i=1}^{N} T_i^n x_i^n
     +\sum_{j=1}^{M} T_j^h x_j^h \\
    &-\sum_{k=1}^{K}\sum_{r=1}^{m_k} c(v_{r-1}^k,v_r^k).
\end{aligned}
\]
Requests are optional. Serving no requests is feasible and gives benefit $0$.
Travel cost includes the leg from a taxi's starting point to its first stop;
there is no return-to-start leg.

\section{Feasibility constraints}

\begin{enumerate}
    \item \textbf{Complete service and assignment.} A served request has one
    pickup and one drop-off, performed in that order by the same taxi. A request
    may instead be left entirely unserved.

    \item \textbf{Time windows.} Passenger pickups and both endpoints of every
    parcel request have time windows. Arrival at such a point must be no later
    than the window's latest time. If arrival precedes its earliest time, the
    taxi waits before starting service. Passenger drop-offs have no time window
    and service begins immediately upon arrival. The stated service duration
    applies at both endpoints of a request.

    \item \textbf{Direct passenger trip.} If taxi $k$ serves passenger $i$,
    $D_i^n$ immediately follows $P_i^n$ in $R_k$. No other request point may
    be visited between them, so the taxi carries at most one passenger at a
    time.

    \item \textbf{Parcel capacity.} Picking up parcel $j$ increases the taxi's
    parcel load by $w_j$; delivering it decreases the load by $w_j$. The load
    must never exceed $Q_k$ for taxi $k$. Passengers do not consume parcel
    capacity, so parcels may remain onboard during a passenger trip.

    \item \textbf{Route timing.} Each taxi starts at $O_k$ at time $0$. Travel
    from consecutive points $u$ and $v$ takes $t(u,v)$, and service at a point
    takes the request's service duration. A taxi may finish at its last service
    point without returning to $O_k$.
\end{enumerate}

\end{document}
